\chapter{Slopes of Compact Hecke Operators}
\label{chap:slopes}

We can now combine the two halves of the thesis.  Chapter~\ref{chap:compact} attaches to a
compact operator its Fredholm determinant and identifies the reciprocal roots with the
eigenvalues; Chapter~\ref{chap:np} reads the valuations of the roots of a power series off the
slopes of its Newton polygon.  Composing:

\begin{center}
  \emph{the slopes of the Newton polygon of $\det(1 - Tu)$ are the valuations of the eigenvalues
  of $u$.}
\end{center}

So to say something about the eigenvalues of a Hecke operator it suffices to bound the
coefficients of its determinant.  This chapter does that in general, and the next specialises to
the case where the bound is attained.

\section{From row decay to coefficient bounds}

Let $u$ be a compactoid operator whose matrix decays along a row weight $w$, meaning $\lVert
u_{ji} \rVert \leq \sigma^{w_j}$.  A minor of $u$ picks entries from $\lvert S \rvert$ distinct
rows, so by the Hadamard bound its norm is at most $\sigma$ to the power of the weight sum over
$S$; the coefficient $c_n$ is an ultrametric sum of $n$-element minors, hence bounded by the
worst such.  The content is therefore combinatorial: how small can $\sum_{j \in S} w_j$ be for
$\lvert S \rvert = n$?

\begin{lemma}
  \label{lem:minor-bound}
  \uses{lem:hadamard, def:iscompactoid}
  \lean{TateFredholm.norm_minor_le_pow_sum}
  \leanok
  If $\lVert u_{ji}\rVert \leq \sigma^{w_j}$ for all $i, j$ then $\lVert \mathrm{minor}(u, S)
  \rVert \leq \sigma^{\sum_{j \in S} w_j}$ for every finite $S$.
\end{lemma}

\begin{theorem}[The slope bound]
  \label{thm:slope-bound}
  \uses{lem:minor-bound, def:minor}
  \lean{TateFredholm.norm_charCoeff_le_pow}
  \leanok
  Let $0 \leq \sigma < 1$, let $u$ be compactoid with $\lVert u_{ji}\rVert \leq \sigma^{w_j}$,
  and let $f : \N \to \N$ satisfy $f(n) \leq \sum_{j \in S} w_j$ for every $S$ with $\lvert S
  \rvert = n$.  Then
  \[
    \lVert c_n(u) \rVert \leq \sigma^{f(n)} .
  \]
\end{theorem}

\begin{proof}
  \leanok
  \uses{lem:minor-bound}
  $c_n$ is $\pm$ the sum of the $n$-element minors, and an ultrametric infinite sum is bounded by
  the supremum of its terms; each term is bounded by Lemma~\ref{lem:minor-bound} and then by
  $\sigma^{f(n)}$ since $\sigma \leq 1$ and $f(n)$ is the smaller exponent.
\end{proof}

Two combinatorial inputs cover the cases we need.

\begin{lemma}
  \label{lem:choose-two}
  \lean{TateFredholm.choose_two_le_sum, TateFredholm.choose_two_lt_sum_of_ne_range}
  \leanok
  Among $n$-element subsets of $\N$, the sum $\sum_{i \in S} i$ is minimised uniquely by $S =
  \{0, 1, \dots, n-1\}$, with value $\binom{n}{2}$; every other $n$-element subset has sum at
  least $\binom{n}{2} + 1$.
\end{lemma}

\begin{lemma}
  \label{lem:block-weight}
  \lean{TateFredholm.sum_div_le_sum_block}
  \leanok
  On $\iota \times \N$ with $\lvert \iota \rvert = d$ and weight $(i, m) \mapsto m$, an
  $n$-element subset has weight sum at least $\sum_{k < n} \lfloor k/d \rfloor$: each value of
  $m$ can be used at most $d$ times.
\end{lemma}

\begin{proof}
  \leanok
  Peel off an element $M$ of maximal second coordinate.  Since $S \subseteq \iota \times \{0,
  \dots, M_2\}$, which has $d(M_2+1)$ elements, we get $n \leq d(M_2+1)$ and hence $\lfloor
  (n-1)/d\rfloor \leq M_2$.  Induct on the remaining $n-1$ elements.
\end{proof}

\begin{corollary}
  \label{cor:slope-bound-block}
  \uses{thm:slope-bound, lem:choose-two, lem:block-weight}
  \lean{TateFredholm.norm_charCoeff_le_pow_choose_two, TateFredholm.norm_charCoeff_le_pow_block}
  \leanok
  On $c(\N, K)$ with $\lVert u_{ji}\rVert \leq \sigma^j$ we get $\lVert c_n \rVert \leq
  \sigma^{\binom{n}{2}}$; on the block model $c(\iota \times \N, K)$ we get $\lVert c_n \rVert
  \leq \sigma^{\sum_{k<n} \lfloor k/\lvert\iota\rvert\rfloor}$.
\end{corollary}

\section{The slope theorem at a general weight}

Applying this to a Hecke operator is now immediate, because
Theorem~\ref{thm:qmf-compact} produced exactly the required row decay.

\begin{proposition}
  \label{prop:qmf-row-decay}
  \uses{thm:qmf-compact, def:heckeblock}
  \lean{QMF.Weight.norm_matrixCoeff_heckeBlock_le, QMF.Weight.norm_matrixCoeff_heckeBlockOp_le}
  \leanok
  With $\lVert \det\theta(\eta)\rVert \leq \sigma$, $\rho \leq \sigma$, and a twisting character
  of norm at most $1$, the block operator satisfies $\lVert \mathrm{matrixCoeff}_{(i,m),(j,r)}
  \rVert \leq \sigma^{m}$.
\end{proposition}

\begin{theorem}
  \label{thm:qmf-slope-bound}
  \uses{prop:qmf-row-decay, cor:slope-bound-block, def:heckecharpowerseries}
  \lean{QMF.Weight.norm_charCoeff_heckeCharPowerSeries_le}
  \leanok
  The coefficients of $\det(1 - T[U\eta U])$ satisfy
  \[
    \lVert c_n \rVert \leq \sigma^{\sum_{k < n}\lfloor k/\lvert \iota\rvert\rfloor}.
  \]
\end{theorem}

Translating a coefficient bound into a polygon statement is what
Definition~\ref{def:ofslopes} was built for.

\begin{definition}
  \label{def:blockslopes}
  \lean{QMF.Weight.blockSlopes, QMF.Weight.monotone_blockSlopes}
  \leanok
  The \emph{block slope sequence} is $s_k = \lfloor k/d \rfloor \cdot (-\log \sigma)$, where $d =
  \lvert \iota \rvert$ is the size of the class set: each of the values $0, -\log\sigma,
  -2\log\sigma, \dots$ repeated $d$ times, once per block.  It is monotone since $\sigma < 1$.
\end{definition}

\begin{theorem}[The slope theorem at a general analytic weight]
  \label{thm:qmf-slope-theorem}
  \uses{thm:qmf-slope-bound, def:blockslopes, def:ofslopes, def:isnewtonpolygonof}
  \lean{QMF.Weight.isBelow_newtonPolygon_heckeCharPowerSeries}
  \leanok
  The Newton polygon of $\det(1 - T[U\eta U])$ lies on or above the polygon anchored at the
  origin whose $k$-th unit slope is $\lfloor k/\lvert\iota\rvert\rfloor \cdot (-\log \sigma)$.
\end{theorem}

\begin{proof}
  \leanok
  \uses{thm:ofslopes, thm:np-exists, thm:qmf-slope-bound}
  The determinant has constant coefficient $1$, so its polygon is anchored at the origin, and all
  its coefficients have norm at most $1$, which is an affine floor and hence admissibility.  The
  comparison polygon is $\mathrm{ofSlopes}$ of the block slope sequence, whose height at $k$ is
  $\sum_{i<k} s_i$; the required inequality height $\leq -\log\lVert c_k\rVert$ is
  Theorem~\ref{thm:qmf-slope-bound} after taking logarithms.  Now apply the maximality clause of
  Definition~\ref{def:isnewtonpolygonof}.
\end{proof}

In words: on a class set of size $d$, the slopes of $U_\varpi$ at weight $\kappa$ are at least
$0, \dots, 0, -\log\sigma, \dots$, each value repeated $d$ times.  This is the general form of
[Jacobs, Theorem 2.12]; the theorem itself asserts equality, and equality needs a hypothesis.

\section{When the bound is attained}

\begin{theorem}[{[Jacobs, Theorem 2.12]}]
  \label{thm:unit-minors}
  \uses{thm:slope-bound, lem:choose-two}
  \lean{JacobsSlash.norm_charCoeff_of_unit_minors}
  \leanok
  Let $u$ on $c(\N,K)$ satisfy $\lVert u_{ji}\rVert \leq \lVert \varpi \rVert^j$ and suppose the
  rescaled matrix $N_{ji} = \varpi^{-j} u_{ji}$ has every top-left minor of norm $1$.  Then
  \[
    \lVert c_m(u) \rVert = \lVert \varpi \rVert^{\binom{m}{2}} \qquad \text{exactly}.
  \]
\end{theorem}

\begin{proof}
  \leanok
  \uses{lem:choose-two, lem:minor-bound}
  The principal block $S = \{0,\dots,m-1\}$ contributes exactly $\lVert\varpi\rVert^{\binom
  m2}$: factoring $\varpi^j$ out of the $j$-th row leaves the unit-minor matrix.  Every other
  $m$-element block contributes strictly less, since its row sum exceeds $\binom m2$ by
  Lemma~\ref{lem:choose-two}.  An ultrametric sum with one strictly dominant term has the norm of
  that term.
\end{proof}

\begin{corollary}[The slope reading]
  \label{cor:slope-reading}
  \uses{thm:unit-minors, thm:ofslopes, cor:ofslopes-transport}
  \lean{JacobsSlash.unitSlope_newtonPolygon₀OfPowerSeries_charPowerSeries}
  \leanok
  Under the hypotheses of Theorem~\ref{thm:unit-minors} the Newton polygon of $\det(1-Tu)$, in
  the normalisation $v(\varpi) = 1$, has unit slopes $0, 1, 2, 3, \dots$.  Consequently the
  eigenvalues of $u$ have valuations $0, 1, 2, \dots$: an arithmetic progression.
\end{corollary}

\begin{proof}
  \leanok
  \uses{thm:ofslopes}
  The valuations $v(c_m) = \binom m2$ are the partial sums of $0,1,2,\dots$, so
  Theorem~\ref{thm:ofslopes} identifies the polygon and
  Corollary~\ref{cor:ofslopes-transport} reads its slopes.  Passing to eigenvalues is
  Theorem~\ref{thm:np-slope-reading} together with Theorem~\ref{thm:serre12}.
\end{proof}

There is a second case, with slopes $\tfrac12, \tfrac32, \tfrac52, \dots$, which occurs when the
field is ramified over $\Qp$ so that half-integers are values of the valuation.  That is what
happens for Jacobs's $U_3$, and it is the subject of the next chapter.
